% Auto-generated by scripts/06_emit_structure.py
% Source: scans 194–199
% Section id: appendix/5
% Phase 3 deliverable. Footnotes (Phase 4), citations (Phase 5),
% and Wade-Giles → Pinyin (Phase 6) are not yet applied here.

\chapter[Relative Chronology]{Relative Chronology: The Periodicity of Divination Topics and Idioms}
\label{app:5}

% source: scan 194, printed 177
% intro-echo-dropped: Relative Chronology:
% intro-echo-dropped: The Periodicity of Divination
% intro-echo-dropped: Topics and Idioms

I.
Introduction
\appendixsectionlabel{sec:appendices-app05:1}{1}
The general evolution of pyromantic practice in the historical period was briefly
described in \ref{sec:chapters-ch04:4.3.1.12} (sec. 4.3.1.12). Changes in the topics and idioms divined can greatly assist our attempts
to periodize particular inscriptions. For the most part, such changes cannot be generalized; they
can only be catalogued. The specific cases listed below, taken from major areas of Shang life, will
alert the reader to a few of these changes and to the ways in which others may be identified as
the occasion arises.
2. Religion
\appendixsectionlabel{sec:appendices-app05:2}{2}
Divinations about sacrifice and the ancestors evolved in a variety of ways. For example,
divinations about ancestral curses were rare after period I.\footnote[1]{Subcharges (\ref{sec:chapters-ch03:3.7.4.2} (sec. 3.7.4.2)) of the form ``It is father (mother, brother, grandfather, grandmother) Chia (Yi, Ping, etc.) who is (is not) cursing (the king)'' (S82.1-83.1) are, with a few possible exceptions, all from period I.} By period V (if not considerably earlier
in some instances), divinations about sacrifice to the various nature powers, such as Ho and
Yüeh, had been discontinued;\footnote[2]{See table 26, nos. 20, 21. Sacrifices to nonancestral powers may have been discontinued as early as period IIb. Whether or not they were restored in periods IIIb and IV (as argued, for example, by \pinyinterm{yan-yiping} [1961], p. 526) depends on the resolution of the Old School-New School controversy (see ch. 2, n. 18).} no divinations were recorded about human sacrifice or about
lao (or, ``penned cattle or sheep'') sacrifice offered to the first twenty-five ancestors in the
sacrificial schedule (that is, the ancestors from Shang Chia to \pinyinterm{xiaoxin}) nor were any
propitiatory or defensive yü rituals divined.\footnote[3]{These conclusions are based on the inscriptions containing ancestral titles listed at S511.1-533.2.}\footnote[4]{See the inscriptions listed at S52.3-56.2.}
Divinations about offering the five major rituals to the ta-tsung ancestors according to a
regular schedule first appeared in period IIb, but it was still necessary, on occasion, to determine
if performing the ritual according to schedule would be acceptable to the powers. The tentative
nature of the sacrificial schedule at this time is demonstrated by the following inscription:
% source: scan 195, printed 178
(1)
(2)
\inscriptionsection{INSCRIPTION}

\begin{inscriptionblock}
\inscriptionref{\pinyinterm{kengyin} divination}

\inscriptionlabel{Preface} Crack-making on \pinyinterm{kengyin} (day 27), [Hsing] divined:
\inscriptionlabel{Charge} ``On the next day, \pinyinterm{xinmao} (day 28), perform hsieh ritual to \pinyinterm{zu-xin} (K13 1); no curse.''
\inscriptionlabel{Postface} In the ninth month.
\end{inscriptionblock}


\inscriptionsection{INSCRIPTION}

\begin{inscriptionblock}
\inscriptionref{ping-shen divination}

\inscriptionlabel{Preface} Crack-making on ping-shen (day 33), Hsing divined:
\inscriptionlabel{Charge} ``On the next day, ting-yu (day 34), perform hsieh ritual to Tsu Ting (K15 1); no curse.''
\inscriptionlabel{Postface} [In the tenth month].\footnote[5]{Jianshou 5.4 (S127.3); also reproduced as \pinyinterm{xubian} 1.18.5. I follow Shima's transcription for ``tenth month,'' which I do not see in the rubbings. If, as Takashima (1976) has proposed, it is safe to deduce subordination from context, the presence of two verbs in the charge suggests a conditional construction: ``If, on the next day, hsin-mao, we perform hsieh-ritual to Tsu Hsin, there will be no curse.'' But cf. n. 6.}
\end{inscriptionblock}


\inscriptionsection{INSCRIPTION}

\begin{inscriptionblock}
\inscriptionref{\pinyinterm{jisi} divination}

(1)
\inscriptionlabel{Preface} [Crack-making] on \pinyinterm{jisi} (day 6), [Hsing] divined:
\inscriptionlabel{Charge} The king [entertains]; performs the chui ritual; [no fault].
\inscriptionlabel{Postface} In the [seventh month].

(2)
\inscriptionlabel{Preface} Crack-making on ping-shen (day 33), Hsing divined:
\inscriptionlabel{Charge} The king entertains Pu Ping (K4 1); performs the hsieh ritual; no fault.
\inscriptionlabel{Postface} In the eighth month.

(3)
\inscriptionlabel{Preface} Crack-making on hsin-ch'ou (day 38), Hsing divined:
\inscriptionlabel{Charge} The king entertains Ta Chia (K3 9)'s consort, Pi Hsin; performs the hsieh ritual; no fault.
\inscriptionlabel{Postface} In the eighth month.

(4)
\inscriptionlabel{Preface} Crack-making on hsin-ch'ou (day 38), Hsing divined:
\inscriptionlabel{Charge} The king entertains; performs the chui ritual.

(5)
\inscriptionlabel{Preface} Crack-making on jen-yin (day 39), Hsing divined:
\inscriptionlabel{Charge} The king entertains Ta Keng (K5 11)'s consort, Pi Jen; performs the hsieh ritual; no fault.

(6)
\inscriptionlabel{Preface} Crack-making on jen-yin (day 39), Hsing divined:
\inscriptionlabel{Charge} The king entertains; performs the chui ritual; no fault.
\inscriptionlabel{Postface} In the eighth month.

(7)
\inscriptionlabel{Preface} Crack-making on jen-tzu (day 49), Hsing divined:
\inscriptionlabel{Charge} The king entertains Ta Wu (K7 1)'s consort, Pi Jen; performs the hsieh ritual; no fault.\footnote[6]{Zhuixin 304. For a good gloss to this inscription and the sacrificial terms, see Ikeda (1964), 1.2.3. It is possible that the charge was conditional: ``If the king entertains Pu Ping and performs the hsieh-ritual, there will be no fault'' (cf. n. 5). But since the divination was performed on the day of the sacrifice (a sacrifice to Pu Ping on ping-shen), it is also possible that the charge was descriptive and hortatory: ``The king is entertaining Pu Ping and performing the hsieh-ritual, and there will be no fault'' (cf. ch. 2, n. 28).}
\end{inscriptionblock}

% source: scan 196, printed 179
rituals (the chui ritual does not belong to the regular five) intrude in what is otherwise an orderly
sequence of rituals addressed to three consorts-recorded in the order of their husband's accession
and place in the sacrificial schedule-who receive cult on the cyclical day corresponding to their
stem name.
Some period II divinations about the five rituals were fully regular, and entire oracle bones
were reserved for this one topic (fig. 18); these inscriptions may thus be regarded as forerunners
of period V divinations about the five rituals that all manifest this regularity unfailingly. By the
reigns of \pinyinterm{di-yi} and \pinyinterm{di-xin}, the ritual schedule had become automatic; the offering of the ritual,
its recipient, and its timing were no longer in question; divinations merely stated that the offering
would be made or was being made on schedule and expressed in formal incantation the general
wish that no fault or misfortune would occur.''\footnote[7]{E.g., Houbian 1.2.12 and other inscriptions listed at S494.2. The yu of the phrase wang yu appears to combine the senses of fault and misfortune.} extraneous divinations never intruded on the bones
and shells reserved for the five-ritual divinations of period V (\ref{sec:chapters-ch04:4.3.1.5} (cf. sec. 4.3.1.5)).
Other chronological variations, based on more technical aspects of Shang divinations about
ritual, may also be established.\footnote[8]{For example, period V charges about the yung, yi, and hsieh rituals invariably add the graph for ``day'' (jih, but always written ; see table 26, no. 10b) after the ritual name; this is found in no other period. See the inscriptions listed at S493.4-494.2; 250.4-251.2; 128.1-2. Houbian 1.3.12 (S128.1) is a typical example: ``(Preface:) Crack-making on chi-hai (day 36), divined: (Charge:) `The king entertains Tsu Ting's (K15 21) consort, Pi Chi; performs the hsieh-day ritual; no fault.' '' The days of these three rituals were used as calendrical markers in late Shang bronze inscriptions; for examples see Akatsuka [1955], pp. 37, 38, 39, 40, together with the ten other cases listed on p. 5 of his index; see too Matsumaru [1970], pp. 65-66, nos. 2-7. Period V divinations may also be distinguished from those of IIb by the occurrence of two or more rituals recorded in the postface (e.g., Jinzhang 518 [D], translated in \ref{sec:chapters-ch04:4.3.1.7} (sec. 4.3.1.7)); inscriptions of this form do not appear in periods I to IV.}
3.
Political Activity
\appendixsectionlabel{sec:appendices-app05:3}{3}
Divinations concerned with political activity and the king's secular interests also reveal
marked periodicity. To cite but a few examples: references to the royal women whose names were
prefixed with fu generally appear only in period I inscriptions and those of the Royal Family
group.\footnote[9]{See the inscriptions listed at S436.2-437.1; cf. Tung (1965), pp. 91-92. Shima (1958), pp. 452-454, lists eighty-six cases of the combination fu-mou; cf. \pinyinterm{zhou-hongxiang-tight} (1970/71), p. 360. Chou's view (pp. 351, 354, 356) that fu-mou also appeared in IV or V is not supported by the evidence he cites. On the status of the royal women, see ch. 1, n. 70.} Certain groups of officers, such as wo shih ``our envoys,'' are recorded only in period
I inscriptions.\footnote[10]{See the inscriptions listed at S421.1.} The names of diviners, even when recorded in the charges and not used as diviners'
names, are frequently limited to the same periods as those of the diviners themselves.\footnote[11]{For example, when Que (S414.3), Cheng (S100.2), Xuan (S314.4-315.2), and Chung (S420.3-4), the names of period I diviners appear in charges, the inscriptions may also be dated to period I. The significance of this will be considered in Studies.} Divinations
seeking to determine if a certain person or group would hsieh wang shih †, ``assist the king's
affairs,'' occur mainly in period I, with a few possible cases in periods II to IV.\footnote[12]{See the inscriptions listed at S124.4-125.2; the possible III + IV exceptions are Jianshou 46.3; \pinyinterm{jingjin} 4777; Nanbei, ``Ming'' 472 (D). The phrase hsieh wo shih +, ``assist our affairs'' (S152.2), occurs only in period I.} Divinations
about tso yi F, ``building a settlement,'' or about tzu yi, ``this settlement,'' occur only in
period I or in Royal Family group inscriptions. Divinations about ta-yi Shang, ``the great
settlement Shang,'' or t'ien-yi Shang, ``the heavenly settlement Shang,'' occur only in
period V.\footnote[13]{See the inscriptions listed at S43.1-3; Shih-to 2.82, to III + IV; the rest are either RFG or, in my opinion, undatable.}

% source: scan 197, printed 180
4. Warfare and the Hunt
\appendixsectionlabel{sec:appendices-app05:4}{4}
In the realm of military affairs, references to Zhi Guo, the Shang general or ally,
appear only in period I,\footnote[14]{See the inscriptions listed at S65.4-67.1; for Zhi Guo, see ch. 3, n. 104.} as do divinations about following or allying with other leaders or
statelets.\footnote[15]{See the inscriptions listed at S20.3-21.4.} Divinations ending with the phrase wo shou yu (=), ``we will receive (ultrahuman) assistance,'' appear only in period I.\footnote[16]{See the inscriptions listed at S351.3-352.1; cf. n. 32.} Divinations about the Shang king raising men rarely
appear after period I.\footnote[17]{See the inscriptions listed at S1.1-3 (but cf. ch. 3, n. 38). I would propose period II datings for Jinzhang 601 (D); Nanbei, ``Ming'' 188 (D); Buci, 131.} Divinations about ``the three hundred archers'' or ``the many archers''
were recorded only in period I and in Royal Family group inscriptions,\footnote[18]{See the inscriptions listed at S378.1-2.} as was the phrase shih
jen \textasciicircum{}, ``to send men.''\footnote[19]{See the inscriptions listed at S1.3-4. Qianbian 8.2.5 (S1.4) is a possible exception; the small calligraphy indicates a period V date, but the bone is fragmentary, and it is not certain that the phrase shih jen is present.} Divinations about 6 \%, ``taking (captured) Ch'iang,'' generally occur
in periods III+IV; they do not occur in I or II.\footnote[20]{See the inscriptions listed at S503.2. On the meaning of an early form of yi Ẻ (D), see CKWT, pp. 4371-4373.} The enemy statelet of the T'u-fang ± was
divined about only in period I;\footnote[21]{See the inscriptions listed at S171.3-172.2; cf. \pinyinterm{chen-mengjia} (1956), pp. 272-273; Tung (1965), p. 90.} the Kung(?)-fang appear in period I and occasionally in
period II;\footnote[22]{See the inscriptions listed at S129.2-132.1; cf. Ch'eng Meng-chia (1956), p. 273; Tung (1965), p. 90.} references to the Chou appear mainly in period I and occasionally in periods II to IV;\footnote[23]{See the inscriptions listed at S299.3-300.1. Shima (1958), p. 409, cites Tongzuan, ``Pieh'' 2. VII.5 and Yizhu go as period II inscriptions referring to Chou. He lists seven period IVb Shiduo 2.82 to III + IV; the rest are either RFG or, in my opinion, undatable. It is a curious fact that though the Chou overthrew the Shang, they are recorded in no period V inscription now available to us.} references to the Yü-fang occur only in period V and are all concerned with \pinyinterm{di-xin}'s campaigns against that statelet.\footnote[24]{See the inscriptions listed at S387.3; cf. \pinyinterm{chen-mengjia} (1956), pp. 309-310. Ch'en, pp. 269-310, gives a period-by-period account of the various statelets that appear in the inscriptions. On the ways in which period Va and Vb inscriptions may or may not be distinguished, see Tung (1945), pt. 1, ch. 3, p. 2a; \pinyinterm{qu-wanli} (1948), p. 225; \pinyinterm{chen-mengjia} (1955), p. 59.}
In the related activity of hunting, the verb varied with time. Shou was commonly used in
period I and more rarely in periods II to IV; it was no longer used in V. T'ien, ``to take the
field,'' hence ``to hunt,'' was used rarely in I but was common in periods II to V.\footnote[25]{See the inscriptions listed at S445.4-447.1; 289.3-298.2. According to Tung (1965), p. 98, t'ien was used in periods II, III, IVa, and V, and shou was used in I and IVb. But I know of no inscription containing shou which can be dated unequivocally to IVb, and I suspect that the inscriptions (unspecified) which he so assigns should be dated rather to the RFG or period III + IV. It seems that t'ien (on the meanings of this word, see \ref{sec:chapters-ch03:3.6.3} (sec. 3.6.3)) and shou were not exact synonyms, for in III + IV they were used jointly in four inscriptions (S446.1; e.g., 甲編 1656; \pinyinterm{jingjin} 4418). The meaning of this compound t'ien-shou, ``to take the field and hunt (?),'' eventually may have come to be assumed by t'ien alone.} The compound
% source: scan 198, printed 181
chu huo ``pursue and catch,'' was used only in period I, as was the verification phrase yün
huo, ``(we) really caught.''\footnote[26]{See the inscriptions listed at S230.2-3; phrases wang hsing 往省, ``go to inspect,'' 489.3-4, and wang t'ien 往田, ``go to hunt'' (S76.4-77.4); the fact that wang hsing was not used after period I suggests that this topic came to be incorporated in the hsing-t'ien, ``inspect and hunt,'' or t'ien-hsing, ``hunt and inspect,'' divinations of III + IV.} References to the king shooting game with bow and arrow
appear only in periods III+IV.\footnote[27]{See the inscriptions listed at S378.4-379.1. Tung (1965), p. 93 (note that Tung [1964], pp. 82-84, consistently misprints \pinyinterm{wu-ding} for \pinyinterm{wu-yi}), dated all these inscriptions to period IVa. This precision was perhaps based upon Tung's desire to find the various classical references (which he cites) to king \pinyinposs{wu-yi} hunting prowess confirmed by the period IVa oracle inscriptions. I find no firm epigraphic or calligraphic justification for a IVa date; \pinyinterm{qu-wanli}, in fact, dates the relevant 甲編 inscriptions to period III.} Similarly, divinations about t'ien-hsing ``hunting and
inspecting,'' in a particular region-which document the political role of the royal hunts--have
only been found for period III + IV.\footnote[28]{See the inscriptions listed at S103.2-4. The first three inscriptions are dated to period I; I take the first two (Buci 203 and Tieyun 114.4) to be an abbreviated or fragmented version of the formula found in the third (Qianbian 5.26.1): hsing wo t'ien, ``inspect our fields.'' Other scholars have taken the period III + IV phrase hsing t'ien to mean ``inspect the hunt'' (e.g., 甲編 1651, kǎoshì, p. 218; 人文 1996, shakubun, p. 504). This is possible, but it fails to account for the fact that the words are frequently reversed, t'ien hsing.} The regions in which the king hunted also varied by period.
Divinations about hunting at Chi, for example, appear only in III+IV, at Sang in III
to V.\footnote[29]{See the inscriptions listed at S469.1; 191.2-4.}
5.
Weather
\appendixsectionlabel{sec:appendices-app05:5}{5}
The weather was a major concern of the Shang. Divinations about the king ``encountering great rain'' (kou ta yü) or ``encountering great wind'' (kou ta feng), which
would affect his hunts and travels, appear mainly in periods II to V.\footnote[30]{See the inscriptions listed at S169.3; 234.1-3; 455-3; 457.2. \pinyinterm{yan-yiping} (1961), p. 536, gives various other rules about the idiom and the context of rain divinations, but I question their accuracy.} Divinations about whether
or not there would be sun appear in periods I to IV and in the inscriptions of the Royal Family
group, but never in period V.\footnote[31]{See the inscriptions listed at S494.4-496.3.} Divinations about ``receiving harvest'' were made throughout
the historical period, but with certain variations in the formula that provide useful dating criteria.\footnote[32]{Charges containing the phrases wo shou nien, ``we will receive harvest'' (fig. 8), were recorded only in period I (S195.2-196.1). In periods III to V, shou nien was replaced by the analogous phrase shou ho (S192.4-193.1). For related epigraphic changes, see table 26, no. 4.3.}
One of the disaster words---thought to refer to drought or dearth---appears only in period I inscriptions.\footnote[33]{The graph (S38.1-3) has been variously interpreted as chin, han, etc. See Shima (1958), p. 192; CKWT, pp. 4013-4018; Mickel (1976), pp. 158-168.}
6.
Other Criteria
\appendixsectionlabel{sec:appendices-app05:6}{6}
An almost endless variety of miscellaneous idioms and graphs reveal periodicities that
may assist us. Words for recording the time of day, the word for consort, disaster formulas, and
% source: scan 199, printed 182
so on are all grist for the mill of the periodizer.\footnote[34]{Time words: The word t'ou, used either to indicate the period between two cyclical days or as a sacrifice term (see ch. 2, n. 79), appears only in period I and RFG inscriptions (S402.3-403.2). The time phrases ta-ts'ai and hsiao-ts'ai, referring approximately to 8 A.M. and 6 P.M. respectively (Tung [1962a], pp. 411-412), appear only in period I and RFG inscriptions (S188.2-3; 497.1), as does the word (=tse ?), ``afternoon'' (S39.1-2). With the exception of its occurrence in the phrase (= lai ch'un [?] 來春 [?]), ``the coming season (?),'' where it also appears in period II, the time word and its variants appear only in period I (S187.2-188.1; cf. Ch'en Mengchia [1956], pp. 227-228; on the meaning of this word, see ch. 3, n. 85). Consort: The graph used in divinations about sacrifice to refer to the royal consorts (e.g., Zhuixin 304 [fig. 7], translated in appendix \ref{app:5}, sec. \ref{sec:appendices-app05:2} (appendix 5, sec. 2)), appears in various epigraphic forms from period IIb to V, but never in period I (S34.4-35.4). Disaster formulas: In the routine ten-day divinations followed by sacrifice dates in the postface (see sec. 4.3.1.7), the king is recorded as the subject of the charge only in period V; in periods II through IV, and occasionally in V (e.g., Qianbian 2.14.1+4.28.1; \pinyinterm{xubian} 1.51.2 [both S168.2]), the subject is always implied. Thus we see a shift in the position of the graph for ``king,'' the earlier formula being wang *tieng \pinyintext{xun} wang huo, ``the king divined: `In the (next) ten days there will be no disaster,' '' the period V formula being *tieng wang \pinyintext{xun} wang huo, ``divined: `the king in the (next) ten days will have no disaster' '' (S166.2-168.2).} %34
 It would be pointless to provide further examples
of this sort here. The relevant criteria may be developed as the occasion arises by consulting the
entries in \pinyinterm{shima-short}'s \booktitle{Sōrui} (\ref{sec:chapters-ch03:3.3.2} (sec. 3.3.2)), by applying the primary and secondary criteria developed in
chapter 4, and by consulting the opinions of other scholars (\ref{sec:chapters-ch04:4.1.1} (sec. 4.1.1)). No date is entirely firm
unless established by the primary criterion of an ancestral title. Many dates offered by scholars,
in fact, depend upon nothing stronger than stylistic or formulaic analogies which may always be
questioned and, if necessary, rejected. Needless to say, future discovery and publication of fresh
stores of oracle bones is likely to alter our view of how the inscriptions evolved.\footnote[35]{Our understanding of period IV in particular, for which we have so few firmly dated, attested inscriptions (see ch. 4, n. 62), may be considerably revised. The majority of the over 4,800 inscribed bone and shell fragments excavated south of \pinyinterm{xiaotun} in 1973 is said to come from periods IVa, IIIb, and, less numerous, IVb (KK [1975.1], p. 39); it will be interesting to see the basis for these dates and the information to be gained from the inscriptions.} %35
